% =========================================================================
% Computationally Efficient Estimation of High-Dimensional Exponential
% Family Models via Tilted Cell Moments and Alternating Projections
%
% Manuscript prepared for submission to a computational statistics journal
% (formatted in the style of the Journal of Computational and Graphical
% Statistics).  Reproduction code: paper/stats/code/ in the accompanying
% repository.
% =========================================================================
\documentclass[11pt]{article}

\usepackage[T1]{fontenc}
\usepackage{lmodern}
\usepackage[utf8]{inputenc}
\usepackage{amsmath,amssymb,amsthm}
\usepackage{bm}
\usepackage{booktabs}
\usepackage{microtype}
\usepackage[margin=1.15in]{geometry}
\usepackage[round,authoryear]{natbib}
\usepackage[colorlinks=true,linkcolor=blue,citecolor=blue,urlcolor=blue]{hyperref}
\usepackage{algorithm}
\usepackage{algpseudocode}
\usepackage{setspace}
\onehalfspacing

\newtheorem{theorem}{Theorem}
\newtheorem{proposition}{Proposition}
\newtheorem{lemma}{Lemma}
\newtheorem{corollary}{Corollary}
\newtheorem{assumption}{Assumption}
\newtheorem{definition}{Definition}
\newtheorem{remark}{Remark}
\newtheorem{example}{Example}

\newcommand{\R}{\mathbb{R}}
\newcommand{\E}{\mathbb{E}}
\newcommand{\Prob}{\mathbb{P}}
\newcommand{\Var}{\operatorname{Var}}
\newcommand{\Cov}{\operatorname{Cov}}
\newcommand{\tr}{\operatorname{tr}}
\newcommand{\diag}{\operatorname{diag}}
\newcommand{\dd}{\,\mathrm{d}}
\newcommand{\convp}{\stackrel{p}{\longrightarrow}}
\newcommand{\convd}{\stackrel{d}{\longrightarrow}}
\newcommand{\bbeta}{\bm{\beta}}
\newcommand{\bgamma}{\bm{\gamma}}
\newcommand{\bdelta}{\bm{\delta}}
\newcommand{\btheta}{\bm{\theta}}
\newcommand{\balpha}{\bm{\alpha}}
\newcommand{\bx}{\mathbf{x}}
\newcommand{\bw}{\mathbf{w}}
\newcommand{\ba}{\mathbf{a}}
\newcommand{\bu}{\mathbf{u}}
\newcommand{\bv}{\mathbf{v}}
\newcommand{\bs}{\mathbf{s}}
\newcommand{\bc}{\mathbf{c}}
\newcommand{\bC}{\mathbf{C}}
\newcommand{\bQ}{\mathbf{Q}}
\newcommand{\bH}{\mathbf{H}}
\newcommand{\bF}{\mathbf{F}}
\newcommand{\bU}{\mathbf{U}}
\newcommand{\bV}{\mathbf{V}}
\newcommand{\bW}{\mathbf{W}}
\newcommand{\bX}{\mathbf{X}}
\newcommand{\bB}{\mathbf{B}}
\newcommand{\bD}{\mathbf{D}}
\newcommand{\bI}{\mathbf{I}}
\newcommand{\bP}{\mathbf{P}}
\newcommand{\bzero}{\mathbf{0}}

\renewcommand{\thetheorem}{S\arabic{theorem}}
\renewcommand{\theproposition}{S\arabic{proposition}}
\renewcommand{\thelemma}{S\arabic{lemma}}
\renewcommand{\thecorollary}{S\arabic{corollary}}
\renewcommand{\theassumption}{S\arabic{assumption}}
\renewcommand{\thedefinition}{S\arabic{definition}}
\renewcommand{\theexample}{S\arabic{example}}
\renewcommand{\theremark}{S\arabic{remark}}
\renewcommand{\thetable}{S\arabic{table}}
\renewcommand{\thefigure}{S\arabic{figure}}
\renewcommand{\theequation}{S\arabic{equation}}
\renewcommand{\thesection}{S\arabic{section}}
\renewcommand{\thealgorithm}{S\arabic{algorithm}}

\title{Online Supplement to ``Tilted Cell Moments: A One-Pass Exact
Newton Method for Generalized Linear Models at Scale''}

\author{Andreas Ljungstr\"om\\
\small Swedish Institute for Social Research, Stockholm University}

\date{July 2026}

\begin{document}

\maketitle

\noindent This supplement develops the companion theory and additional
methodology for the framework of the main paper: global convergence and
local rates for the nonlinear alternating-projections absorption of
high-dimensional fixed effects (Sections~\ref{sec:hdfe}, with the
$G \ge 3$ rate theory in Section~\ref{sec:multiG}); asymptotic theory in
fixed, diverging, and proportional-dimension regimes
(Section~\ref{sec:asymptotics}); further extensions---families beyond
Poisson at high dimension, penalized estimation, and Firth's
bias-reducing adjustment (Section~\ref{sec:extensions}); simulation
studies validating each result (Section~\ref{sec:supp-sims}); and all
proofs.  Notation follows the main paper; numbered objects here carry
the prefix S.

\section*{Notation recalled}

Observations $i = 1, \dots, N$ carry responses $y_i$ from an exponential
dispersion family with cumulant $b$, individual-level covariates
$\bx_i \in \R^p$, cell-level covariates $\bw_{j(i)} \in \R^k$ constant
within the cells $j = 1, \dots, J$ of the categorical
cross-classification, and offsets $o_i$; $\bbeta = (\bgamma, \bdelta)$
collects the individual- and cell-level coefficients, and for canonical
links $\omega_i = b''(\eta_i)$.  The tilted cell moments of main-paper
Theorem~1 are
\begin{equation}\label{eq:s-moments}
M_j^{(0)} = \sum_{i \in j} \omega_i, \qquad
\mathbf{m}_j = \sum_{i \in j} \omega_i \bx_i, \qquad
\bQ = \sum_{i=1}^N \omega_i \bx_i \bx_i' ,
\end{equation}
from which the exact score and information are assembled at
$O(N(p^2 + p)) + O(J(k^2 + kp))$ per Newton or Fisher-scoring iteration.
Algorithm~S1 below denotes the absorbed-factor scheme
analyzed here: cyclic closed-form profile updates of the absorbed
effects alternated with damped tilted-moment Newton steps for $\bbeta$.

% =========================================================================
\section{High-dimensional fixed effects and nonlinear alternating
projections}\label{sec:hdfe}
% =========================================================================

\subsection{Profile formulation}\label{sec:profile}

Beyond the cell design, applications routinely require one or more
fixed-effect factors whose level counts preclude even cell-level
treatment: firms, hospitals, zip-code--year interactions.  Augment the
linear predictor with $G$ absorbed factors,
\begin{equation}\label{eq:linpred-fe}
\eta_i = \bgamma'\bx_i + \bdelta'\bw_{j(i)} + \sum_{g=1}^{G}
  \alpha^{(g)}_{g(i)} + o_i ,
\end{equation}
where factor $g$ has $J_g$ levels with level map $g(i)$ and effect vector
$\balpha^{(g)} \in \R^{J_g}$, and $J_1 + \dots + J_G$ may be of order
$N$.  The full parameter is
$\btheta = (\bbeta, \balpha)$ with $\bbeta = (\bgamma, \bdelta)$.  The
design has a known null space $\mathcal{K} \subset \R^{q + \sum_g J_g}$
(at least the $G$ dimensions of shifts between the factors' levels and
the intercept, more if the factor graph is disconnected); all statements
below are on the quotient $\R^{q+\sum J_g}/\mathcal{K}$, realized in
practice by a normalization.  The estimator of interest is the MLE of
$\bbeta$ with $\balpha$ concentrated out: the profile likelihood
$\ell_p(\bbeta) = \max_{\balpha} \ell(\bbeta, \balpha)$.

For the Poisson family the inner maximization has a closed-form
coordinate solution: given all other parameters, the optimal effect of
level $a$ of factor $g$ is
\begin{equation}\label{eq:closed-form}
\alpha^{(g)}_a
 = \log \frac{\sum_{i : g(i) = a} y_i}
             {\sum_{i : g(i) = a} \exp\!\left(\eta_i - \alpha^{(g)}_a\right)},
\end{equation}
the log of the ratio of observed to expected-at-zero-effect events, which
is the multiplicative-update rule used by \texttt{FENmlm}/\texttt{fixest}
\citep{berge2018} and, in IRLS-linearized form, by \texttt{ppmlhdfe}
\citep{correia2020}.  Cycling \eqref{eq:closed-form} over factors is a
nonlinear block Gauss--Seidel scheme; combined with tilted-moment Newton
steps for $\bbeta$ it yields Algorithm~\ref{alg:absorb}.

\begin{algorithm}[t]
\caption{Tilted-moment Newton with absorbed factors (nonlinear block
Gauss--Seidel)}\label{alg:absorb}
\begin{algorithmic}[1]
\State Initialize $\bbeta_0$, $\balpha_0 = \bzero$; drop observations in
absorbed groups with no events (Section~\ref{sec:separation}).
\Repeat
  \Repeat \Comment{inner loop: exact block maximization over $\balpha$}
    \For{$g = 1, \dots, G$}
      \State update $\balpha^{(g)}$ by the closed-form profile rule
        \eqref{eq:closed-form}
    \EndFor
  \Until{$\max_g \|\Delta\balpha^{(g)}\|_\infty < \tau_{\mathrm{in}}$}
  \State accumulate tilted moments \eqref{eq:s-moments} at
    $(\bbeta_t, \balpha_t)$ in one microdata pass
  \State $\bbeta_{t+1} \gets \bbeta_t + \lambda_t\, \bF^{-1}\bU$
    with $\lambda_t \in \{1, \tfrac12, \tfrac14, \dots\}$ chosen by
    backtracking (Armijo) on $\ell(\cdot, \hat\balpha(\cdot))$
\Until{$\max_\ell |\Delta\beta_\ell| / (1 + |\beta_\ell|) < \tau$}
\end{algorithmic}
\end{algorithm}

Algorithm~\ref{alg:absorb} is, up to the tilted-moment acceleration of
the $\bbeta$-step, the algorithm used across the applied HDFE--GLM
software ecosystem.  Its convergence status splits cleanly in two.  The
\emph{inner} loop is the method of alternating projections, whose
convergence is classical---von Neumann for two subspaces, generalized
by \citet{halperin1962}, with the Kaczmarz sweep \citep{kaczmarz1937}
as the coordinate case---and which \citet{gaure2013} introduced into
fixed-effects estimation on exactly this foundation;
\citet{correia2020} prove convergence of the linear inner solves they
employ, and \citet{deutschhundal1997} supply rates.  For the
\emph{combined} nonlinear cycle---profile updates of the absorbed
effects alternated with Newton/scoring steps on the GLM
likelihood---the literature offers simulation evidence
\citep{guimaraes2010,berge2018,correia2020} but no unified formal
theorem.  The next two subsections supply that missing statement:
Theorem~\ref{thm:global} (main-paper Theorem~2) proves global
convergence of the combined inner--outer iteration, and
Theorems~\ref{thm:rate}--\ref{thm:multiG} characterize its rate.

\subsection{Global convergence}\label{sec:global}

We impose:

\begin{assumption}\label{ass:canonical}
The link is canonical and $b$ is strictly convex and twice continuously
differentiable on $\Theta = \R$.
\end{assumption}

\begin{assumption}\label{ass:mle-exists}
The MLE exists: $\sup_{\btheta} \ell$ is attained at some
$\btheta^\ast$ (unique modulo $\mathcal{K}$).
\end{assumption}

\begin{assumption}\label{ass:blocks}
(i) Each absorbed group contains at least one observation with $y_i$ in
the interior of the convex hull of the support of the family (for
Poisson and binomial: at least one $y_i > 0$, and for binomial also one
$y_i < n_i$), so that each one-dimensional block likelihood in
\eqref{eq:closed-form} attains a finite maximum; (ii) the $\bbeta$-block
information $\bF_{\bbeta\bbeta}(\btheta)$ is nonsingular on the initial
superlevel set $\{\ell \ge \ell(\btheta_0)\}$.
\end{assumption}

Assumption~\ref{ass:mle-exists} is a data condition---absence of
separation---with a long literature of verifiable characterizations
\citep{haberman1974,wedderburn1976,albert1984,santossilva2010,correia2019};
Section~\ref{sec:separation} discusses enforcement and the Firth
alternative when it fails.  Assumption~\ref{ass:blocks}(i) is enforced
constructively by the drop rule of Algorithm~\ref{alg:absorb}, mirroring
standard practice \citep{correia2020}.

\begin{theorem}[Global convergence of nonlinear alternating
projections; main-paper Theorem~2]\label{thm:global}
Under Assumptions~\ref{ass:canonical}--\ref{ass:blocks}, the iterates
$\btheta_t = (\bbeta_t, \balpha_t)$ of Algorithm~\ref{alg:absorb}
satisfy: (i) $\ell(\btheta_t)$ is nondecreasing and converges to
$\sup \ell = \ell(\btheta^\ast)$; (ii) $[\btheta_t] \to [\btheta^\ast]$
in the quotient $\R^{q+\sum J_g}/\mathcal{K}$; in particular
$\bbeta_t \to \bbeta^\ast$ whenever the normalization places the flat
directions in the absorbed blocks; and (iii) the same conclusions hold if
the $\bbeta$-update is replaced by exact block maximization, and for any
cyclic or essentially-cyclic ordering of the blocks; and (iv)
Assumption~\ref{ass:canonical} may be weakened to the requirement that
the link be twice continuously differentiable with each per-observation
log-likelihood contribution strictly concave in its linear predictor
(a \emph{log-concave link}; \citealp{pratt1981}), provided the
$\bbeta$-step is an Armijo-damped ascent step with a symmetric
preconditioner whose eigenvalues are bounded within $[\underline\lambda,
\bar\lambda] \subset (0,\infty)$ on the initial superlevel
set---Fisher scoring in particular.
\end{theorem}

Part (iv) settles the scope of the guarantee beyond canonical links.
Log-concavity of the per-observation contribution in $\eta_i$ is exactly
the condition under which the proof's three ingredients
survive---concavity of $\ell$ in $\btheta$, unique finite block
maximizers, and sufficient ascent of a preconditioned gradient
step---and it is satisfied well beyond the canonical family:
\citet{pratt1981} characterizes the log-concave links, which include
the gamma family with log link (each contribution
$-\nu\{y_i e^{-\eta_i} + \eta_i\}$ has second derivative
$-\nu y_i e^{-\eta_i} < 0$), the binomial probit, and the binomial
complementary log-log.  For the gamma--log pair the absorbed blocks even
retain a closed-form update,
$\alpha^{(g)}_a = \log\{ m_a^{-1} \sum_{i : g(i)=a} y_i
\exp(-(\eta_i - \alpha^{(g)}_a))\}$ with $m_a$ the group size, so
Algorithm~\ref{alg:absorb} applies unchanged.  For links outside this
class the log likelihood need not be concave in $\btheta$ and the
guarantee weakens to what smooth ascent methods generically provide:
monotone likelihood ascent with stationary limit points, without a
global claim; we state this boundary explicitly since it is where the
folk confidence in the algorithm stops being a theorem.

The proof (Appendix~\ref{app:global}) proceeds by establishing that
strict concavity on the quotient plus attainment forces compact
superlevel sets (Lemma~\ref{lem:levelsets}); monotonicity then confines
the iterates to a compact set, every limit point is shown to be a
blockwise maximizer using the uniqueness of the block updates and the
sufficient-ascent property of the damped Newton step, and for smooth
concave functions blockwise maximizers are global maximizers.  The
argument is in the tradition of block-coordinate methods
\citep{bertsekas1999,tseng2001}; what requires care is (a) the quotient
structure (the likelihood is not coercive in the natural parametrization,
only modulo $\mathcal{K}$), (b) the inexact $\bbeta$-block (a damped
Newton step, not an exact maximizer), and (c) dispensing with any
compactness or boundedness assumption on the parameter space, which would
be false here and is where informal appeals to standard results break
down.

Theorem~\ref{thm:global} has a practical corollary worth stating: since
every accumulation point of the algorithm is \emph{the} MLE, the familiar
diagnostics of the linear-projection literature (sensitivity of estimates
to tightened tolerances, discrepancies across implementations) reflect
finite-iteration error only, not multiplicity; and step-halving on the
profile objective---re-solving the inner loop at every trial point, as
Algorithm~\ref{alg:absorb} specifies---is guaranteed to terminate at each
outer iteration.

\subsection{Local rate: Friedrichs angles and a bipartite spectral
gap}\label{sec:rate}

Convergence speed is governed by the geometry of the absorbed subspaces.
Let $\Omega^\ast = \diag(\omega_i^\ast)$ collect the information weights
at the optimum, let
$V_g = \operatorname{span}\{\text{indicator columns of factor } g\}
\subset \R^N$, and equip $\R^N$ with the inner product
$\langle u, v \rangle_{\Omega^\ast} = u'\Omega^\ast v$.  Recall that the
\emph{Friedrichs angle} $\theta_F \in (0, \pi/2]$ between closed
subspaces $V_1, V_2$ with intersection $V_0$ has cosine
\[
c_F = \sup\left\{ \frac{\langle u, v\rangle_{\Omega^\ast}}
  {\|u\|_{\Omega^\ast} \|v\|_{\Omega^\ast}} :
  u \in V_1 \cap V_0^{\perp},\;
  v \in V_2 \cap V_0^{\perp},\; u, v \ne \bzero \right\} < 1 .
\]

\begin{theorem}[Local linear convergence]\label{thm:rate}
Let Assumptions~\ref{ass:canonical}--\ref{ass:blocks} hold and let $b$ be
three times continuously differentiable.  (i) The inner iteration of
Algorithm~\ref{alg:absorb} at fixed $\bbeta$ (the map
$T : \balpha \mapsto \balpha^{+}$ of one Gauss--Seidel sweep of exact
block maximizations) is continuously differentiable in a neighborhood of
its fixed point $\hat\balpha(\bbeta)$, with Jacobian equal to the block
Gauss--Seidel iteration matrix of the $\balpha$-block information
$\bF_{\balpha\balpha}$ at the fixed point; consequently the inner
iteration converges locally $R$-linearly with rate
$\rho_{\mathrm{GS}} = \rho\big( (\bD + \mathbf{L})^{-1} \mathbf{L}' \big)
< 1$ on the quotient, where $\bF_{\balpha\balpha} = \bD + \mathbf{L} +
\mathbf{L}'$ is the block decomposition into diagonal and strictly
lower/upper block-triangular parts.
(ii) For $G = 2$ absorbed factors,
\[
\rho_{\mathrm{GS}} = c_F^2 ,
\]
the squared Friedrichs cosine between $V_1$ and $V_2$ in the
$\Omega^\ast$ inner product; the same rate governs the weighted
alternating-projections demeaning of any working vector, by the theorem
of \citet{aronszajn1950} and \citet{kayalar1988}.
(iii) The outer iteration (full cycles of Algorithm~\ref{alg:absorb} with
exact $\bbeta$-maximization) converges locally $R$-linearly with rate at
most $\rho\big( \text{block GS operator of } \bF(\btheta^\ast) \text{ in
the ordering } (\balpha^{(1)}, \dots, \balpha^{(G)}, \bbeta) \big) < 1$.
\end{theorem}

\begin{corollary}[Bipartite spectral gap]\label{cor:bipartite}
For $G = 2$, define the weighted biadjacency matrix
$\bB \in \R^{J_1 \times J_2}$ by
$B_{ab} = \sum_{i : g_1(i) = a,\, g_2(i) = b} \omega_i^\ast$, and let
$\bD_1 = \diag(\bB \mathbf{1})$, $\bD_2 = \diag(\bB' \mathbf{1})$.  Then
\[
c_F = \sigma_2\!\left( \bD_1^{-1/2} \bB\, \bD_2^{-1/2} \right),
\]
the second-largest singular value of the normalized biadjacency
(transition) kernel.  Equivalently, $1 - c_F^2$ is the spectral gap of
the random walk on the weighted bipartite graph of the two factors, and
$c_F < 1$ if and only if that graph is connected---the classical
identification condition for two-way fixed effects
\citep{abowd1999,gaure2013}.
\end{corollary}

Corollary~\ref{cor:bipartite} turns the folklore statement
``correlated fixed effects make demeaning slow'' into a computable
quantity: before iterating at all, $\sigma_2$ of a $J_1 \times J_2$
sparse matrix (a few Lanczos iterations) predicts the per-sweep
contraction factor exactly.  It also gives sharp intuition for the
pathologies: nearly disconnected mobility graphs---workers who rarely
switch firms, patients who rarely switch hospitals---have
$\sigma_2 \uparrow 1$ and predictably slow demeaning, while balanced
crossed designs have small $\sigma_2$ and converge in a handful of
sweeps.  Section~\ref{sec:sim-rate} verifies the prediction to three
decimals across two orders of magnitude of $1 - c_F^2$.

\subsection{Three or more absorbed factors}\label{sec:multiG}

Administrative designs routinely absorb three or more factors---worker,
firm, and region--year, say---and for $G \ge 3$ the linear
alternating-projections literature offers norm bounds but no clean angle
characterization of the asymptotic rate \citep{deutsch2001,badea2012}.
The following theorem resolves the practical question for our setting:
the exact local rate for any $G$ is a computable spectral quantity
determined by the \emph{pairwise} weighted contingency structure of the
factors alone, so the bipartite intuition of
Corollary~\ref{cor:bipartite} does not break down at $G \ge 3$---it
generalizes, with the pair graph replaced by a $G$-partite sweep
operator.  For each pair $g \ne h$ define the weighted cross-tabulation
$\bB_{gh} \in \R^{J_g \times J_h}$,
$(\bB_{gh})_{ab} = \sum_{i : g(i)=a,\, h(i)=b} \omega_i^\ast$, and the
level totals $\bD_g = \diag(\bB_{gh}\mathbf{1})$ (independent of $h$).

\begin{theorem}[Exact rate and certificates for $G \ge 2$ factors]
\label{thm:multiG}
Let the assumptions of Theorem~\ref{thm:rate} hold and let
$\mathcal{F} = \bF_{\balpha\balpha}(\btheta^\ast)$ denote the absorbed-block
information on the level space $\R^{J_1 + \dots + J_G}$.
\begin{enumerate}
\item[(i)] \emph{(Pairwise sufficiency.)}  $\mathcal{F}$ has diagonal
blocks $\bD_g$ and off-diagonal blocks $\bB_{gh}$; hence the local rate
$\rho_{\mathrm{GS}}$ of Theorem~\ref{thm:rate}(i) is a function of the
pairwise weighted cross-tabulations only, and can be evaluated without
access to the microdata once the $G(G-1)/2$ tables are accumulated (one
$O(N)$ pass).
\item[(ii)] \emph{(Exact computable rate.)}  Write $\mathcal{F} =
\mathcal{D} + \mathcal{L} + \mathcal{L}'$ with $\mathcal{D} =
\diag(\bD_1, \dots, \bD_G)$ and $\mathcal{L}$ strictly block lower
triangular, and let $\mathbf{M} = -(\mathcal{D} +
\mathcal{L})^{-1}\mathcal{L}'$ be the sweep operator.  Then the
eigenspace of $\mathbf{M}$ for the eigenvalue $1$ is exactly
$\operatorname{null}(\mathcal{F})$ (the flat shift directions), this
eigenvalue is semisimple, and
\[
\rho_{\mathrm{GS}}
 = \max\big\{ |\lambda| : \lambda \in \operatorname{spec}(\mathbf{M}),\;
   \lambda \ne 1 \big\} < 1
 \iff \operatorname{null}(\mathcal{F}) = \text{flats}
 \iff \text{the multi-way design is identified}.
\]
Because the eigenvalue $1$ is semisimple, successive differences of the
sweep iteration decay at exactly $\rho_{\mathrm{GS}}$, so the rate is
computable by power iteration on the level space at
$O\big(\sum_{g<h} \operatorname{nnz}(\bB_{gh})\big)$ per step.
\item[(iii)] \emph{(A priori certificate.)}  With $c_g$ the Friedrichs
cosine, in the $\Omega^\ast$ inner product, between $V_g$ and
$V_{g+1} + \dots + V_G$ for $g = 1, \dots, G-1$,
\[
\rho_{\mathrm{GS}}
 \;\le\; \big\| \text{one demeaning sweep} \big\|_{\Omega^\ast}
 \;\le\; \Big( 1 - \prod_{g=1}^{G-1} \big(1 - c_g^2\big) \Big)^{1/2},
\]
each $c_g$ again being a level-space quantity (the largest
non-unit generalized singular value between factor $g$ and the pooled
design of the later factors); $c_g < 1$ for all $g$ if and only if the
design is identified.
\item[(iv)] For $G = 2$, (ii) reduces to $\rho_{\mathrm{GS}} =
c_F^2 = \sigma_2^2$ of Corollary~\ref{cor:bipartite}, and the bound in
(iii) to $c_F$.
\end{enumerate}
\end{theorem}

Three comments.  First, part (ii) is the operative solution for
practice: the exact rate is available \emph{before iterating on the
microdata}, from cross-tabulations that cost one linear pass, by power
iteration on vectors of length $\sum_g J_g$; the reference
implementation exposes this as \texttt{gs\_spectral\_rate}, and
Section~\ref{sec:sim-rate} shows it matching observed contraction
factors to three decimals for $G = 3$ exactly as for $G = 2$.  Second,
part (iii) gives the interpretable geometry: each successive factor
slows the sweep according to its angle to the span of the factors after
it, the product-of-sines form being the multi-subspace bound of
\citet{smith1977} and \citet{deutschhundal1997} transported through the
complement duality; the bound is attained-in-order for $G = 2$ but is
conservative for $G \ge 3$ (Table~\ref{tab:sim2g3}), which is precisely
why (ii), not an angle formula, is the right object to compute.  The
sharp one-sweep \emph{energy-norm} contraction for general $G$ is
characterized abstractly by the identity of \citet{xu2002} for
multiplicative subspace-correction methods---the demeaning sweep is
exactly such a method---but that identity's constant is itself a
variational quantity no easier than (ii), reinforcing the same
conclusion.  Third, the identification equivalence in (ii) generalizes
the connectedness criterion of \citet{abowd1999} and \citet{gaure2013}:
for $G \ge 3$, pairwise connectedness of all factor pairs is necessary
but not sufficient for $\rho_{\mathrm{GS}} < 1$ on the quotient; the
spectral condition is exactly right, and the power iteration returns
$\rho_{\mathrm{GS}} = 1$ (to tolerance) when identification fails,
making it a cheap diagnostic as well as a rate predictor.

\subsection{Separation}\label{sec:separation}

Assumption~\ref{ass:mle-exists} fails in the maximum-likelihood sense
when a hyperplane in covariate space separates events from non-events:
components of $\hat\btheta$ diverge and the supremum is not attained
\citep{albert1984,santossilva2010}.  In the absorbed blocks the failure
mode is transparent---a group with no events has
$\hat\alpha_a = -\infty$ by \eqref{eq:closed-form}---and
Algorithm~\ref{alg:absorb} drops such groups a priori, after which
Assumption~\ref{ass:blocks}(i) holds by construction; this mirrors
\citet{correia2019}, who treat detection in full generality.  Separation
induced by the $\bbeta$-block cannot be dropped away and surfaces as
non-convergence.  Section~\ref{sec:firth} develops the principled remedy
---Firth's adjusted score inside the tilted accumulation---and
Section~\ref{sec:sim-separation} quantifies its behavior.

% =========================================================================
\section{Asymptotic theory}\label{sec:asymptotics}
% =========================================================================

Because Theorem~1 of the main paper makes the tilted-moment estimator
\emph{equal} to the MLE, its asymptotic analysis is the analysis of the
MLE in the relevant regime.  We give results in three regimes,
in increasing order of dimensionality.  Throughout,
$\btheta_0$ denotes the truth, expectations are under the model, and the
triangular-array framework lets all design quantities depend on $N$.

\subsection{Fixed dimension}\label{sec:classical}

\begin{assumption}\label{ass:classical}
(i) $q = p + k$ is fixed; covariates and offsets are uniformly bounded;
(ii) $\lambda_{\min}\big(\bF_N(\btheta_0)\big) \to \infty$ and
$N^{-1}\bF_N(\btheta_0) \to \bF_\infty \succ \bzero$;
(iii) $\btheta_0$ is interior to the parameter space and $b$ is three
times continuously differentiable.
\end{assumption}

\begin{theorem}[Classical asymptotics]\label{thm:classical}
Under Assumptions~\ref{ass:canonical} and \ref{ass:classical}, with
probability tending to one the MLE $\hat\btheta_N$ exists, is unique, and
satisfies
$\sqrt{N}\, (\hat\btheta_N - \btheta_0) \convd
\mathcal{N}\big(\bzero,\, \phi\,\bF_\infty^{-1}\big)$; the estimator is
asymptotically efficient, attaining the Cram\'er--Rao bound.  If the
conditional mean is correctly specified but the distribution is not, the
same limit holds with sandwich variance
$\bF_\infty^{-1} \bm{\Sigma} \bF_\infty^{-1}$, consistently estimated by
the cell-accumulated sandwich of Section~4.4 of the main paper.
\end{theorem}

The proof (Appendix~\ref{app:classical}) is a specialization of
\citet{fahrmeir1985} plus the pseudo-ML results of
\citet{gourieroux1984}; we include it because the triangular-array
statement with offsets and the exact form of the cell-accumulated
sandwich are used repeatedly later.

\subsection{Diverging dimension: $q_N^3/N \to 0$}\label{sec:diverging}

Register applications place $k_N = \dim(\bdelta)$ in the hundreds or
thousands---period $\times$ region $\times$ education cells---so the
fixed-$q$ regime is a poor description.  Let $q_N = p + k_N \to \infty$,
write $\bF_N = \bF_N(\btheta_0)$, and define the \emph{leverage} of
observation $i$ as $h_i = \omega_i\, \ba_i' \bF_N^{-1} \ba_i$ (so
$\sum_i h_i = q_N$).

\begin{assumption}\label{ass:diverging}
(i) Covariates and offsets are uniformly bounded, and the true linear
predictors $\{\eta_{0i}\}$ lie in a fixed compact set;
(ii) \emph{balanced leverage}: $\max_i h_i \le C\, q_N / N$ for a fixed
constant $C$;
(iii) $q_N^3 / N \to 0$;
(iv) $b$ is three times continuously differentiable with $b''$ bounded
away from $0$ and $\infty$, and $b'''$ bounded, on compact sets.
\end{assumption}

Condition (ii) is the natural design condition: for a pure cell design
with balanced cells it holds with $C$ close to $1$, and it fails exactly
when some cell is drastically undersized relative to the average---the
same configurations for which practitioners distrust the asymptotics.

\begin{theorem}[Diverging dimension]\label{thm:diverging}
Under Assumptions~\ref{ass:canonical} and \ref{ass:diverging}, with
probability tending to one the MLE exists and satisfies
\[
\big\| \bF_N^{1/2} (\hat\btheta_N - \btheta_0) \big\|
  = O_p\big(\sqrt{q_N}\big) ,
\]
and for any fixed-dimensional matrix of contrasts $\bC_N' \in
\R^{d \times q_N}$ with $\bC_N' \bF_N^{-1} \bC_N \to \bm{\Sigma}_C \succ
\bzero$,
\[
\big(\bC_N' \bF_N^{-1} \bC_N\big)^{-1/2} \bC_N'
  (\hat\btheta_N - \btheta_0) \convd \mathcal{N}(\bzero, \phi\,\bI_d) .
\]
The plug-in variance estimator $\bC_N'\bF_N(\hat\btheta_N)^{-1}\bC_N$ is
consistent in relative terms.
\end{theorem}

The proof (Appendix~\ref{app:diverging}) is self-contained and short by
the standards of this literature, because concavity in the canonical
parametrization substitutes for the empirical-process machinery needed in
general $M$-estimation \citep{portnoy1988,heshao2000,wang2011}: a
quadratic-minorization argument in the intrinsic norm delivers the rate,
and a third-order Taylor control with the balanced-leverage condition
delivers asymptotic linearity with remainder $o_p(1)$ exactly when
$q_N^3/N \to 0$.

\subsection{The proportional regime $J/N \to c$: Poisson profile
likelihood as conditional likelihood}\label{sec:proportional}

The regime that motivates this paper's title question is more extreme:
one absorbed effect per cell with cells of \emph{bounded} size, so the
parameter count is a positive fraction of the sample size and
Theorem~\ref{thm:diverging} cannot apply ($q_N \asymp N$).  This is the
incidental-parameters regime of \citet{neyman1948}; for smooth nonlinear
panel models the profile MLE of the common parameter is typically
inconsistent or biased at order $1/m$ for cell size $m$
\citep{lancaster2000,fernandezval2016}.  The Poisson family with log
link is a distinguished exception, for a structural reason worth
displaying: profiling and conditioning coincide.

Consider cells $j = 1, \dots, J$ of sizes $m_j \le \bar m < \infty$, one
free effect $\alpha_j$ per cell, individual-level covariates
$\bx_{ij}$, offsets $o_{ij}$, and
$y_{ij} \sim \mathrm{Poisson}\big(\exp(\alpha_j + \bgamma'\bx_{ij} +
o_{ij})\big)$ independently.  Let $D_j = \sum_i y_{ij}$.

\begin{lemma}[Profile equals conditional]\label{lem:conditional}
For every $\bgamma$ with all $D_j$-positive cells,
\[
\ell_p(\bgamma)
 \;=\; \max_{\balpha} \ell(\bgamma, \balpha)
 \;=\; \sum_{j : D_j > 0}
   \log \Prob\big( y_{j1}, \dots, y_{jm_j} \,\big|\, D_j ; \bgamma \big)
   \;+\; \text{terms not involving } \bgamma ,
\]
where the conditional law of $(y_{j1},\dots,y_{jm_j})$ given $D_j$ is
multinomial with $D_j$ trials and probabilities
$\pi_{ij}(\bgamma) = e^{\bgamma'\bx_{ij} + o_{ij}} / \sum_{i'}
e^{\bgamma'\bx_{i'j} + o_{i'j}}$.  Moreover $\ell_p$ is concave in
$\bgamma$, and its score and negative Hessian are the Schur complements
of the tilted moments:
\[
\bU_p(\bgamma) = \sum_j \sum_i \big( y_{ij} - D_j \pi_{ij} \big)
  \bx_{ij},
\qquad
\bH_p(\bgamma) = \sum_j D_j\, \Var_{\pi_j}(\bx_j)
 = \bQ - \sum_j \frac{\mathbf{m}_j \mathbf{m}_j'}{M_j^{(0)}} .
\]
\end{lemma}

Lemma~\ref{lem:conditional} is a fixed-effects analogue of the classical
Poisson--multinomial connection; versions of the first display appear in
the econometric panel literature
\citep{andersen1970,hausman1984,blundell2002,wooldridge1999}, and the
Schur-complement identity is what the tilted-moment machinery computes at
$O(Np^2)$ per iteration (function \texttt{profile\_poisson} in the
accompanying code).  Its consequence, however, seems not to have been
developed as an explicit high-dimensional limit theorem with variance
theory, which we now state.  Conditional on covariates and cell effects,
cells are independent; define the average conditional information
$\bar{\bF}_J(\bgamma_0) = J^{-1} \sum_j \E\big[ D_j
\Var_{\pi_j}(\bx_j) \big]$.

\begin{assumption}\label{ass:proportional}
(i) Cell sizes satisfy $2 \le m_j \le \bar m$; $J/N \to c \in (0, 1)$;
(ii) covariates and offsets are uniformly bounded and $\sup_j \alpha_j
\le C$, so cell means are bounded above (they may be arbitrarily small);
(iii) $\bar{\bF}_J(\bgamma_0) \to \bar{\bF} \succ \bzero$ as
$J \to \infty$;
(iv) $\bgamma_0$ is interior to a fixed compact set $\Gamma$.
\end{assumption}

\begin{theorem}[Proportional regime]\label{thm:proportional}
Under Assumption~\ref{ass:proportional}: (i) the profile MLE
$\hat\bgamma_J$ is consistent; (ii)
\[
\sqrt{J}\, \big( \hat\bgamma_J - \bgamma_0 \big)
 \convd \mathcal{N}\big( \bzero,\; \bar{\bF}^{-1} \big) ,
\]
with \emph{no asymptotic bias}, despite the number of nuisance
parameters being a positive fraction of $N$; (iii) both the profile
information estimator $J\,\bH_p(\hat\bgamma_J)^{-1}$ and the cell-clustered
sandwich estimator
$J\,\bH_p^{-1} \big( \tfrac{G}{G-1}\sum_j \bs_j \bs_j' \big) \bH_p^{-1}$,
with $\bs_j$ the per-cell profile score, are consistent for the
asymptotic variance in (ii)---the latter remaining consistent under
within-cell dependence misspecification with correctly specified
conditional means; and (iv) $\hat\bgamma_J$ attains the semiparametric
efficiency bound for the model in which the $\alpha_j$ are unrestricted
nuisance parameters, the conditional score being the efficient score.
\end{theorem}

Three comments.  First, the mechanism is special to the Poisson--log
pair: the profile score is \emph{exactly} the conditional score, hence
exactly unbiased for every $J$ and $m$; the incidental-parameter bias
that plagues, e.g., the fixed-effects logit (where conditional ML exists
but differs from profile ML, and profile ML has $O(1/m)$ bias
\citep{andersen1970,fernandezval2016}) is identically zero here.
Section~\ref{sec:sim-logit} exhibits the logit failure mode and its
conditional-likelihood repair quantitatively, and
Section~\ref{sec:families} discusses how both the conditional
likelihood \citep{chamberlain1980} and the analytic bias corrections of
\citet{hahn2004} and \citet{fernandezval2016} slot into the
cell-collapsed computational structure.  This is
the likelihood-theoretic content of the folklore that ``Poisson FE is
immune to the incidental-parameters problem''
\citep{lancaster2000,wooldridge1999}, sharpened to a $J/N \to c$ limit
theorem with variance estimation.  Second, the theorem's variance is the
\emph{conditional} information, which the naive full-information plug-in
would overstate; the tilted-moment machinery produces the correct Schur
complement at no extra cost.  Third, event-free cells contribute no
information and are dropped; this is the profile-likelihood counterpart
of the separation drop rule and is accounted for in the theory by summing
over $D_j > 0$.

Section~\ref{sec:sim-highdim} probes the theorem at $c = 1/2$---one
fixed effect per two observations---where nominal $95\%$ intervals from
both estimators in (iii) cover at $94$--$96\%$.

% =========================================================================
\section{Further extensions}\label{sec:extensions}
% =========================================================================

\subsection{Families beyond Poisson}\label{sec:families}

The framework of Section~2 of the main paper was stated for a generic
exponential dispersion family; we record the two most used instances.

\emph{Logistic regression.}  With $b(\theta) = \log(1 + e^{\theta})$,
$\mu_i = (1+e^{-\eta_i})^{-1}$ and $\omega_i = \mu_i(1-\mu_i)$: the tilt
weights each within-cell observation by its Bernoulli variance, so the
tilted measure concentrates on observations near the decision boundary.
All of Theorem~1 of the main paper, Corollary~1 of the main paper, and
Theorems~\ref{thm:global}--\ref{thm:multiG} apply verbatim (the logit
link is canonical).  Unlike the Poisson case the tilt depends on
$\bdelta$ through $\omega_i$, so the within-cell moments must be
re-accumulated each iteration---which the single pass does anyway.  The
proportional-regime theory of Section~\ref{sec:proportional} does
\emph{not} transfer: profiling and conditioning differ for the logit,
and profile ML carries $O(1/m)$ incidental-parameter bias---for matched
pairs ($m = 2$) the classical limit is $2\gamma_0$
\citep{andersen1970}---which Section~\ref{sec:sim-logit} reproduces
numerically.  Two repairs are compatible with the cell-collapsed
computational structure.  \emph{Conditional ML}
\citep{chamberlain1980}: conditioning on cell totals gives per-cell
denominators that are elementary symmetric functions of
$\{e^{\bgamma'\bx_{ij}}\}_{i \in j}$; these obey the standard
Newton--Girard recursion and are computable in $O(m_j D_j)$ operations
per cell---polynomial, not combinatorial---so for the bounded cell sizes
of the proportional regime the conditional score and information are one
more $O(N)$ cell pass (the reference implementation includes an exact
conditional logit along these lines, used in
Section~\ref{sec:sim-logit}).  \emph{Analytic bias correction}
\citep{hahn2004,fernandezval2016}: the correction terms are averages of
within-cell weighted moments of $(\mu_i, \omega_i,
\partial\omega_i/\partial\eta)$ and leverage-type quantities, all of
which are cell-decomposable in exactly the sense of
Theorem~1 of the main paper and Proposition~\ref{prop:firth}; there is no
computational barrier to evaluating them inside the tilted pass, and we
view a packaged implementation as worthwhile software work rather than
an open methodological problem.

\emph{Gamma regression with log link.}  Here the log link is not
canonical, Proposition~2 of the main paper applies, and the Fisher
weights are identically one: the ``tilted'' moments are ordinary
exposure-weighted cell moments, constant across iterations.  One pass
suffices for the whole optimization, and only the working residual sums
$R_j$ are refreshed---a family where the acceleration is at its most
extreme.

\subsection{Penalized estimation}\label{sec:penalized}

In applications with many continuous covariates (genomics: expression
levels alongside batch cells), sparsity-inducing penalties on $\bgamma$
are natural while leaving $\bdelta$ and the absorbed effects
unpenalized.  Consider
\[
\min_{\bgamma, \bdelta}\;
 -\ell(\bgamma, \bdelta) + \lambda_1 \|\bgamma\|_1
 + \tfrac{\lambda_2}{2} \|\bgamma\|_2^2 .
\]

\begin{proposition}[Penalized tilted-moment proximal Newton]
\label{prop:penalized}
At each outer iteration, form the quadratic model of $-\ell$ from the
tilted moments \eqref{eq:s-moments} (one $O(Np^2)$ pass), i.e., the exact
second-order expansion with gradient $-\bU$ and curvature $\bF$.
(i) Ridge ($\lambda_1 = 0$): the update solves the $(p+k)$-dimensional
system with $\lambda_2 \bI$ added to the $\bgamma$-block; all conclusions
of Corollary~1 of the main paper hold.
(ii) Lasso/elastic net: one cycle of coordinate descent on the penalized
quadratic costs $O((p+k)^2)$---independent of $N$---after the pass, so
the inner solver of \citet{friedman2010} runs entirely on cell-collapsed
quantities; the outer iteration is a proximal Newton method and converges
quadratically to the penalized optimum near the solution under the
conditions of \citet{leesun2014}.
(iii) The Karush--Kuhn--Tucker conditions, hence the active set and the
regularization path, are functions of the tilted moments only.
\end{proposition}

The practical consequence: a lasso path for a Poisson model with
$N = 10^7$, $J = 10^3$ cells, and $p = 50$ expression covariates costs
per path point what a dense implementation spends per coordinate.

\subsection{Firth's adjustment inside the tilted pass}\label{sec:firth}

Under separation (Section~\ref{sec:separation}) the MLE does not exist;
Firth's bias-reducing penalty $\ell^\ast(\btheta) = \ell(\btheta) +
\tfrac12 \log\det \bF(\btheta)$ \citep{firth1993} guarantees a finite
maximizer for binomial models \citep{heinze2002} and removes the
$O(N^{-1})$ bias term in general.  For canonical links the adjusted score
has the closed form
$\bU^\ast = \bU + \sum_i h_i\, \frac{b'''(\eta_i)}{2 b''(\eta_i)}\,
\ba_i$ with leverages $h_i = \omega_i \ba_i' \bF^{-1} \ba_i$
\citep{firth1993,kosmidis2009}.

\begin{proposition}[Tilted Firth]\label{prop:firth}
Given the partitioned inverse of $\bF$ (available from the Newton solve),
the leverages decompose as
$h_i = \omega_i \{ \bx_i' \bV_{\gamma\gamma} \bx_i
 + 2\, \bx_i' \bc_{j(i)} + d_{j(i)} \}$ with cell-level precomputations
$\bc_j = \bV_{\gamma\delta} \bw_j$ and
$d_j = \bw_j' \bV_{\delta\delta} \bw_j$ at $O(J(kp + k^2))$; the adjusted
score, and hence the entire Firth quasi-scoring iteration
\citep{kosmidis2009}, is computable at
$O(N(p^2+p)) + O(J(k^2+kp)) + O((p+k)^3)$ per iteration---the same order
as Corollary~1 of the main paper.  For Poisson,
$b'''/2b'' \equiv 1/2$; for logit it is $\tfrac12(1 - 2\mu_i)$.
\end{proposition}

Thus the statistically principled response to separation costs
essentially nothing beyond the unpenalized fit, removing the last reason
to respond to separation by ad hoc data deletion.
Section~\ref{sec:sim-separation} quantifies the finite-sample behavior.

A clarification on software behavior is in order.  Both implementations
detect and drop separated \emph{absorbed} groups automatically before
estimation (Section~\ref{sec:separation}).  Neither switches to the
Firth estimator automatically: separation induced by the $\bbeta$-block
surfaces as flagged non-convergence (the Stata command sets
\texttt{e(converged)} $= 0$; the reference implementation returns
\texttt{converged=False}), and the Firth fit must be requested
explicitly (\texttt{firth=True} in the reference implementation; the
Stata command currently reports the diagnosis and does not expose a
Firth option).  The non-switching is deliberate: the Firth estimator is
a different estimand-defining penalty, not a numerical fallback, and
silently substituting it would change reported quantities without the
user's knowledge.  A principled automatic \emph{detector}---running the
linear-programming certificate of \citet{correia2019} on the cell
aggregates, at $O(J)$ rather than $O(N)$ scale---is a natural software
addition that the tilted representation makes cheap.

% =========================================================================
\section{Simulation studies for the supplement results}
\label{sec:supp-sims}
% =========================================================================

The studies below validate the results of
Sections~\ref{sec:hdfe}--\ref{sec:extensions}; they follow the
computational protocol of main-paper Section~5 (single laptop core,
fixed seeds, scripts \texttt{sim2}--\texttt{sim5} in the code
supplement) and complete in minutes.

\subsection{Convergence rates of alternating projections}
\label{sec:sim-rate}

Study 2 validates Theorem~\ref{thm:rate}(ii) and
Corollary~\ref{cor:bipartite}.  Two factors ($J_1 = 100$, $J_2 = 80$,
$N = 2\times10^5$) are assigned with a mixing parameter $\rho$: with
probability $\rho$ a unit's level of factor 2 is deterministically tied
to its level of factor 1, otherwise uniform---so $\rho$ sweeps the
weighted bipartite graph from expander-like ($\rho = 0$) to nearly
disconnected ($\rho = 0.99$).  At Poisson optimum weights we compute the
predicted per-sweep rate $\sigma_2^2$ from the normalized biadjacency
kernel, then measure the empirical contraction factor of weighted
alternating-projections demeaning of a generic covariate.

\begin{table}[t]
\centering
\caption{Study 2: predicted versus empirical per-sweep contraction of
weighted alternating projections; sweeps to $\|\Delta\|_\infty <
10^{-13}$.}
\label{tab:sim2}
\small
\begin{tabular}{ccccc}
\toprule
$\rho$ & Predicted $c_F^2 = \sigma_2^2$ & Empirical rate & Sweeps \\
\midrule
0.00 & 0.0024 & 0.0019 & 7 \\
0.30 & 0.1897 & 0.1873 & 19 \\
0.60 & 0.5567 & 0.5547 & 48 \\
0.80 & 0.7542 & 0.7534 & 96 \\
0.90 & 0.8896 & 0.8747 & 197 \\
0.95 & 0.9469 & 0.9469 & 468 \\
0.99 & 0.9911 & 0.9910 & 2487 \\
\bottomrule
\end{tabular}
\end{table}

Table~\ref{tab:sim2} shows agreement to the third decimal across the
entire range, including the near-disconnected regime where the rate is
within $10^{-3}$ of $1$ and the sweep count grows as
$(1 - c_F^2)^{-1}$.  The spectral prediction is computable before
iterating (a sparse SVD of the $100 \times 80$ weighted cross-tabulation)
and could be used to warn users, select orderings, or switch to
conjugate-gradient acceleration when the gap is small
\citep{gaure2013,correia2016}.

The second panel of the study probes Theorem~\ref{thm:multiG} with
$G = 3$ factors ($J_1, J_2, J_3 = 100, 80, 60$; factor 2 tied to factor
1 and factor 3 to factor 2, each with probability $\rho$).  For each
$\rho$ we compute the exact level-space spectral rate of
Theorem~\ref{thm:multiG}(ii) (function \texttt{gs\_spectral\_rate},
an eigencomputation on $240$ levels), the Friedrichs product certificate
of Theorem~\ref{thm:multiG}(iii), and the empirical per-sweep
contraction of the three-factor demeaning.

\begin{table}[t]
\centering
\caption{Study 2, $G = 3$ absorbed factors: exact spectral rate of
Theorem~\ref{thm:multiG}(ii) versus empirical per-sweep contraction,
with the product-of-sines norm certificate of
Theorem~\ref{thm:multiG}(iii).}
\label{tab:sim2g3}
\small
\begin{tabular}{ccccc}
\toprule
$\rho$ & Exact $\rho_{\mathrm{GS}}$ & Empirical rate & Norm bound &
Sweeps \\
\midrule
0.00 & 0.0066 & 0.0064 & 0.0766 & 7 \\
0.30 & 0.2520 & 0.2524 & 0.6218 & 23 \\
0.60 & 0.5293 & 0.5317 & 0.8788 & 46 \\
0.80 & 0.7907 & 0.7890 & 0.9763 & 113 \\
0.90 & 0.8802 & 0.8837 & 0.9928 & 225 \\
0.95 & 0.9548 & 0.9552 & 0.9989 & 569 \\
\bottomrule
\end{tabular}
\end{table}

Table~\ref{tab:sim2g3} confirms the two halves of the theorem exactly
as designed: the spectral rate matches the observed contraction to the
third decimal at every $\rho$---the $G = 2$ precision carries over to
$G = 3$ unchanged---while the a priori norm bound is valid everywhere
but conservative, overstating the rate substantially in the
well-conditioned regimes.  The practical protocol this supports: use the
product certificate (computable from angles factor-by-factor) for a
quick feasibility screen, and the $O(\sum J_g)$-dimensional spectral
computation for the exact rate and for the identification diagnostic of
Theorem~\ref{thm:multiG}(ii).

\subsection{The proportional regime}\label{sec:sim-highdim}

Study 3 probes Theorem~\ref{thm:proportional} at its most extreme:
Poisson cells of bounded size $m \in \{2, 5, 20\}$, one fixed effect per
cell, heterogeneous cell effects
$\alpha_j \sim \mathcal{N}(-1.2, 0.8^2)$, a covariate correlated with the
cell effects (so that ignoring the effects biases $\hat\bgamma$), $N \in
\{20{,}000,\ 80{,}000\}$, $1{,}000$ replications.  With $m = 2$ the
design estimates one nuisance parameter per two observations
($c = 1/2$).

\begin{table}[t]
\centering
\caption{Study 3: profile (concentrated) Poisson estimation with
$J/N = 1/m$; $1{,}000$ replications; $\gamma_{01} = 0.3$.  ``Cover'' is
the empirical coverage of nominal $95\%$ intervals from the profile
information (P) and the cell-clustered sandwich (C).}
\label{tab:sim3}
\small
\begin{tabular}{rrcccccc}
\toprule
$N$ & $m$ & $J/N$ & Bias$(\hat\gamma_1)$ & SD$(\hat\gamma_1)$ &
$\overline{\mathrm{SE}}_P$ & Cover$_P$ & Cover$_C$ \\
\midrule
20{,}000 & 2  & 0.50 & $-0.0003$ & 0.0156 & 0.0157 & 0.954 & 0.953 \\
20{,}000 & 5  & 0.20 & $+0.0009$ & 0.0118 & 0.0119 & 0.953 & 0.951 \\
20{,}000 & 20 & 0.05 & $+0.0004$ & 0.0107 & 0.0107 & 0.954 & 0.953 \\
80{,}000 & 2  & 0.50 & $+0.0004$ & 0.0080 & 0.0078 & 0.953 & 0.951 \\
80{,}000 & 5  & 0.20 & $-0.0001$ & 0.0059 & 0.0060 & 0.955 & 0.954 \\
80{,}000 & 20 & 0.05 & $-0.0001$ & 0.0054 & 0.0053 & 0.943 & 0.941 \\
\bottomrule
\end{tabular}
\end{table}

Table~\ref{tab:sim3} confirms every element of the theorem: bias is an
order of magnitude below the Monte Carlo standard deviation even at
$c = 1/2$ and shrinks with $N$ at fixed $c$; the profile-information
standard errors track the true sampling variability to the third decimal;
and both interval constructions cover at $94$--$96\%$ throughout.  For
contrast, the same data analyzed with the cell effects ignored (a single
intercept) biases the coefficient of the effect-correlated covariate by
$+0.28$ (truth $-0.2$) and yields nominal $95\%$ intervals with
\emph{zero} empirical coverage in every configuration---the point being
that the fixed effects are not ignorable here, they are merely free.

\subsection{Separation and the Firth-adjusted tilted pass}
\label{sec:sim-separation}

Study 4 stresses the boundary where Assumption~\ref{ass:mle-exists}
fails.  Logistic model, $J = 10$ cells, a continuous covariate, and a
rare binary exposure ($3\%$ prevalence, true log-odds $1.5$): in small
samples the exposed events can be quasi-separated with appreciable
probability.  We compare the MLE (tilted Newton, non-convergence or
$|\hat\gamma_z| > 10$ flagged as separated) with the Firth-adjusted
tilted iteration of Proposition~\ref{prop:firth}, $500$ replications per
$N$.

\begin{table}[t]
\centering
\caption{Study 4: quasi-separation frequency and estimator behavior;
true $\gamma_z = 1.5$.  MLE summaries condition on non-separated
replications; Firth summaries use all replications.  MAE is median
absolute error.}
\label{tab:sim4}
\small
\begin{tabular}{rcccccc}
\toprule
$N$ & Separated & Firth converged & med$(\hat\gamma_z^{\mathrm{MLE}})$ &
med$(\hat\gamma_z^{\mathrm{Firth}})$ & MAE$_{\mathrm{MLE}}$ &
MAE$_{\mathrm{Firth}}$ \\
\midrule
300  & 0.394 & 1.000 & 1.704 & 1.586 & 0.587 & 0.498 \\
600  & 0.042 & 1.000 & 1.581 & 1.556 & 0.446 & 0.400 \\
1200 & 0.002 & 1.000 & 1.519 & 1.512 & 0.235 & 0.222 \\
2400 & 0.000 & 1.000 & 1.520 & 1.515 & 0.199 & 0.195 \\
\bottomrule
\end{tabular}
\end{table}

At $N = 300$ the MLE fails to exist in $39\%$ of replications; the Firth
iteration converges in all $500$, with smaller median bias
\emph{conditional on the MLE existing at all} and uniformly smaller
median absolute error.  As $N$ grows the two estimators merge, as theory
requires.  The cost of the adjustment was a factor $\sim\!2$ per
iteration (the extra unadjusted pass for current-iterate leverages),
consistent with Proposition~\ref{prop:firth}.

\subsection{The logit failure mode and its conditional repair}
\label{sec:sim-logit}

Study 5 quantifies the boundary drawn in
Sections~\ref{sec:proportional} and~\ref{sec:families}: in the
proportional regime the Poisson profile MLE is exactly unbiased
(Theorem~\ref{thm:proportional}, Study 3), but the logit profile MLE is
not.  Cells of size $m \in \{2, 4, 8\}$, one logit fixed effect per
cell, $\alpha_j \sim \mathcal{N}(0,1)$, a single covariate with
$\gamma_0 = 1$, $N = 8{,}000$, $500$ replications; uninformative cells
(all events or all non-events) are dropped, as both estimators require.
We compare the profile MLE (inner concentration of the cell effects by
vectorized Newton, outer Schur-complement steps---the logit analogue of
\texttt{profile\_poisson}) with the exact conditional MLE of
\citet{chamberlain1980}, computed cellwise as described in
Section~\ref{sec:families}.

\begin{table}[t]
\centering
\caption{Study 5: incidental-parameter bias of the logit profile MLE at
$J/N = 1/m$ versus the conditional MLE; true $\gamma_0 = 1$; $500$
replications.  The $m = 2$ profile mean approaches Andersen's classical
$2\gamma_0$ limit.}
\label{tab:sim5}
\small
\begin{tabular}{rccccc}
\toprule
$m$ & $J/N$ & Mean $\hat\gamma^{\mathrm{profile}}$ &
Bias$^{\mathrm{profile}}$ & Mean $\hat\gamma^{\mathrm{cond}}$ &
Bias$^{\mathrm{cond}}$ \\
\midrule
2 & 0.50  & 1.869 & $+0.869$ & 1.000 & $-0.000$ \\
4 & 0.25  & 1.403 & $+0.403$ & 0.998 & $-0.002$ \\
8 & 0.125 & 1.165 & $+0.165$ & 0.998 & $-0.002$ \\
\bottomrule
\end{tabular}
\end{table}

Table~\ref{tab:sim5} displays the textbook pathology at full strength:
the profile bias is $87\%$ of the parameter at $m = 2$ (approaching the
$2\gamma_0$ limit of \citealp{andersen1970}) and decays at the
predicted $O(1/m)$, while the conditional estimator is unbiased to the
third decimal at every $m$ \emph{at the same cell-collapsed
computational cost}.  Read together with Study 3, the two tables make
the paper's asymptotic message concrete: the proportional regime is not
generically safe (logit) nor generically fatal (Poisson); it is the
profile--conditional coincidence, an algebraic property of the family,
that decides---and when the coincidence fails, the conditional
likelihood is the drop-in repair within the same framework.

\appendix
\renewcommand{\thesection}{S-\Alph{section}}

\section{Proofs for Section~\ref{sec:hdfe} of this supplement}\label{app:global}
% =========================================================================

Throughout this appendix write $D = q + \sum_g J_g$ for the total
parameter dimension, $\ba_i \in \R^{D}$ for the stacked design vector of
observation $i$ (individual-level, cell-level, and absorbed indicator
parts), $\mathcal{K} = \{ \bu : \ba_i'\bu = 0 \ \forall i \}$ for the
null space, and $E = \mathcal{K}^{\perp} \cong \R^{D}/\mathcal{K}$.  By
Assumption~\ref{ass:canonical}, $\ell(\btheta) = \phi^{-1}\sum_i \{y_i
\ba_i'\btheta - b(\ba_i'\btheta + o_i)\}$ (absorbing offsets into $b$'s
argument) is concave on $\R^D$, constant on cosets of $\mathcal{K}$, and
its restriction $\bar\ell$ to $E$ is \emph{strictly} concave: for
$\bu \in E \setminus \{\bzero\}$, some $\ba_i'\bu \ne 0$ and
$-\bu'\nabla^2\ell\,\bu = \phi^{-1}\sum_i b''(\cdot)(\ba_i'\bu)^2 > 0$.

\begin{lemma}[Compact superlevel sets]\label{lem:levelsets}
Let $f : \R^d \to \R$ be strictly concave, continuous, and attain its
maximum at $x^\ast$.  Then every superlevel set
$\{x : f(x) \ge a\}$, $a \le f(x^\ast)$, is compact.
\end{lemma}

\begin{proof}
Superlevel sets of concave functions are closed and convex.  Suppose
$S_a$ is unbounded for some $a$; then $S_a$ contains a ray
$\{x^\ast + t\bv : t \ge 0\}$ for some unit vector $\bv$ (a closed
unbounded convex set containing $x^\ast$ contains a ray from $x^\ast$;
see \citealp[Thm.~8.4]{rockafellar1970}).  Let
$\varphi(t) = f(x^\ast + t\bv)$: strictly concave on $[0, \infty)$,
$\varphi(t) \ge a$ for all $t$, maximized at $t = 0$.  If
$\varphi'(t_0^+) < 0$ for some $t_0$, concavity gives $\varphi(t) \le
\varphi(t_0) + \varphi'(t_0^+)(t - t_0) \to -\infty$, contradicting
$\varphi \ge a$; hence the (strictly decreasing, by strict concavity)
right derivative satisfies $\varphi'(t^+) \ge 0$ for all $t > 0$.  But
$\varphi'(0^+) \le 0$ at the maximum, and strict decrease then forces
$\varphi'(t^+) < 0$ for $t > 0$---a contradiction.
\end{proof}

\subsection{Proof of Theorem~\ref{thm:global}}

All statements are on $E$; representatives are normalized after each
sweep (any fixed measurable normalization, e.g., recentering each
absorbed factor and compensating in the intercept, changes nothing in
$\ell$).

\emph{Monotonicity and confinement.}  Each $\balpha^{(g)}$-update is an
exact maximization over its block (the block objective separates over
the levels of factor $g$; each level's scalar problem $\alpha \mapsto
c_1 \alpha - \sum_{i} b(\alpha + \text{rest}_i)$ is strictly concave
with derivative $\to \pm\infty$ prevented only if the level has an
interior event, which Assumption~\ref{ass:blocks}(i) guarantees; hence a
unique finite maximizer exists).  The $\bbeta$-update is accepted only
under the Armijo test, so $\ell(\btheta_{t+1}) \ge \ell(\btheta_t)$, and
the sequence $\ell(\btheta_t)$ converges to some $L^\ast \le
\ell(\btheta^\ast)$ (Assumption~\ref{ass:mle-exists}).  By
Lemma~\ref{lem:levelsets} applied to $\bar\ell$ on $E$, the normalized
iterates remain in the compact set $\bar S = \{\bar\ell \ge
\ell(\btheta_0)\}$.

\emph{Limit points are stationary.}  Let $\bar\btheta$ be a limit point
of the normalized iterates, along $t_r$.  The per-cycle increment
$\ell(\btheta_{t+1}) - \ell(\btheta_t) \to 0$, and increments are sums
of nonnegative per-block increments, so each block's increment vanishes.
Fix a block $g$ and let $T_g$ denote the map that exactly maximizes over
block $g$ holding the rest fixed; $T_g$ is continuous on $\bar S$ (the
unique maximizer of a strictly concave, jointly continuous objective
depends continuously on the fixed coordinates; Berge's maximum theorem
plus uniqueness).  If block $g$ were not optimal at $\bar\btheta$, then
$\varepsilon := \ell(T_g\bar\btheta) - \ell(\bar\btheta) > 0$; by
continuity of $T_g$ and $\ell$, the block-$g$ increment at
$\btheta_{t_r}$ (with the preceding blocks of the cycle already updated,
which also converge along the subsequence by the same argument applied
inductively in the cycle order) exceeds $\varepsilon/2$ for large $r$
---contradiction.  Hence every absorbed block of $\bar\btheta$ is
blockwise optimal, i.e., $\nabla_{\balpha}\ell(\bar\btheta) = \bzero$.
For the $\bbeta$-block, suppose $\bU_{\bbeta}(\bar\btheta) \ne \bzero$.
On the compact set $\bar S$, the information $\bF$ is continuous and,
by Assumption~\ref{ass:blocks}(ii), $\bF_{\bbeta\bbeta}$ is nonsingular
with eigenvalues in a fixed interval $[\underline\lambda,
\bar\lambda]$, $\underline\lambda>0$; the Newton direction $\mathbf{d} =
\bF_{\bbeta\bbeta}^{-1}\bU_{\bbeta}$ then satisfies $\bU_{\bbeta}'\mathbf{d}
\ge \|\bU_{\bbeta}\|^2/\bar\lambda$, and the standard backtracking bound
guarantees an accepted step length $\lambda_t$ bounded below and an
increment $\ge \sigma \lambda_t \bU_{\bbeta}'\mathbf{d}$ bounded away from
zero in a neighborhood of $\bar\btheta$ \citep[Prop.~1.2.1 and
Sec.~2.7]{bertsekas1999}---again contradicting vanishing increments.
Hence $\nabla\ell(\bar\btheta) = \bzero$.

\emph{Conclusion.}  For a concave differentiable function, a stationary
point is a global maximizer, so $\bar\btheta = \btheta^\ast$ modulo
$\mathcal{K}$ by strict concavity on $E$; since every limit point of the
normalized iterates in the compact set $\bar S$ equals the normalized
$\btheta^\ast$, the whole sequence converges to it, and
$L^\ast = \ell(\btheta^\ast)$, proving (i)--(ii).  Part (iii) follows
because the argument used only: exactness/uniqueness of the absorbed
block maximizations, sufficient ascent of the $\bbeta$-step, and
vanishing per-block increments---all invariant to replacing the
$\bbeta$-step by exact maximization and to essentially-cyclic orderings
(every block updated at least once every $B$ iterations).

For part (iv), inspect where Assumption~\ref{ass:canonical} entered:
(a) concavity of $\ell$ on $\R^D$ and strict concavity on $E$ used only
that each per-observation contribution $\eta \mapsto \ell_i(\eta)$ is
(strictly) concave, which is the log-concave-link hypothesis; (b) the
absorbed block subproblems are separable sums of strictly concave
scalar functions, so unique finite maximizers exist under the analogue
of Assumption~\ref{ass:blocks}(i) (for the gamma--log family every
group maximizer is finite whenever the group is nonempty, since
$y_i > 0$ a.s.); and (c) the sufficient-ascent argument for the
$\bbeta$-step used only that the step direction is
$\mathbf{d} = \mathbf{P}^{-1}\bU_{\bbeta}$ for a symmetric
preconditioner $\mathbf{P}$ with eigenvalues in $[\underline\lambda,
\bar\lambda]$, giving $\bU_{\bbeta}'\mathbf{d} \ge
\|\bU_{\bbeta}\|^2/\bar\lambda$---satisfied by the expected
information of Fisher scoring under
Assumption~\ref{ass:blocks}(ii). \qed

\subsection{Proof of Theorem~\ref{thm:rate}}

\emph{(i)}  Write the inner objective as $f(\balpha) =
\ell(\bbeta, \balpha)$ with fixed $\bbeta$, and let $\hat\balpha$ be its
maximizer (unique after normalization).  The block-$g$ exact update
solves $\nabla_g f(\balpha^{(<g),+}, \alpha^{(g)},
\balpha^{(>g)}) = \bzero$; the block Hessian $\bF_{gg}$ is diagonal with
entries $\sum_{i \in \text{level}} b''(\eta_i) > 0$, so by the implicit
function theorem (using $b \in C^3$) the update map is $C^1$ near the
fixed point with partial derivatives $-\bF_{gg}^{-1}\bF_{g, h}$ with
respect to block $h \neq g$.  Composing the $G$ block updates in cycle
order gives $DT(\hat\balpha) = -(\bD + \mathbf{L})^{-1}\mathbf{L}'$,
where $\bF_{\balpha\balpha} = \bD + \mathbf{L} + \mathbf{L}'$ is the
block diagonal/lower/upper decomposition---the classical Gauss--Seidel
error operator---restricted to the quotient.  Since
$\bF_{\balpha\balpha}$ is symmetric positive definite on the quotient,
$\rho_{\mathrm{GS}} = \rho(DT) < 1$ \citep[Sec.~11.2]{golub2013}, and
Ostrowski's theorem \citep[Sec.~10.1]{ortega1970} yields local $R$-linear
convergence at rate $\rho_{\mathrm{GS}}$.

\emph{(ii)}  Let $G = 2$ with indicator matrices $\mathbf{A}_1 \in
\R^{N\times J_1}$, $\mathbf{A}_2 \in \R^{N \times J_2}$ and $\bF_{gh} =
\mathbf{A}_g'\Omega^\ast\mathbf{A}_h$.  The eigenvalues of the two-block
Gauss--Seidel operator are $\{0\}$ together with the eigenvalues of
$\bF_{22}^{-1}\bF_{21}\bF_{11}^{-1}\bF_{12}$ (eliminate the first block).
In the $\Omega^\ast$ inner product this matrix represents
$\bP_2 \bP_1|_{V_2}$, the compression to $V_2$ of the product of the
orthogonal projections onto $V_1$ and $V_2$; its spectrum on the
complement of $V_1 \cap V_2$ is $\{\cos^2\theta_i\}$ with $\theta_i$ the
principal angles between $V_1 \ominus (V_1\cap V_2)$ and $V_2 \ominus
(V_1 \cap V_2)$, whose supremum-cosine is $c_F$
\citep[Ch.~9]{deutsch2001}.  On the quotient (which removes
$V_1 \cap V_2$, the flat directions), $\rho = c_F^2$.  The identical
rate for the demeaning residuals is the theorem of
\citet{aronszajn1950} and \citet{kayalar1988}:
$\|(\bP_{V_2^\perp}\bP_{V_1^\perp})^n - \bP_{(V_1+V_2)^\perp}
\|_{\Omega^\ast} = c_F^{2n-1}$, using that the Friedrichs angle is
invariant under orthogonal complementation
\citep[Ch.~9]{deutsch2001}.

\emph{(iii)}  Append the $\bbeta$-block: the cycle map with exact
maximizations is again $C^1$ by the implicit function theorem (the
$\bbeta$-block Hessian is nonsingular by
Assumption~\ref{ass:blocks}(ii)), its Jacobian is the block Gauss--Seidel
operator of the full $\bF(\btheta^\ast)$ in the stated ordering, and the
SPD argument of (i) applies on the quotient. \qed

\subsection{Proof of Corollary~\ref{cor:bipartite}}

Represent $u \in V_1$, $v \in V_2$ by their level vectors, $u =
\mathbf{A}_1\tilde u$, $v = \mathbf{A}_2 \tilde v$.  Then
$\langle u, v\rangle_{\Omega^\ast} = \tilde u' \bB \tilde v$,
$\|u\|^2_{\Omega^\ast} = \tilde u' \bD_1 \tilde u$,
$\|v\|^2_{\Omega^\ast} = \tilde v' \bD_2 \tilde v$.  Substituting
$\mathbf{r} = \bD_1^{1/2}\tilde u$, $\mathbf{s} = \bD_2^{1/2}\tilde v$
turns the Friedrichs supremum into the second singular value of
$\mathbf{T} = \bD_1^{-1/2}\bB\bD_2^{-1/2}$: the top singular pair
$(\bD_1^{1/2}\mathbf{1}, \bD_2^{1/2}\mathbf{1})$, with $\sigma_1 = 1$,
spans exactly the image of the intersection $V_1 \cap V_2 \supseteq
\text{constants}$, which the Friedrichs angle excludes.  $\mathbf{T}$ is
the transition kernel of the two-step random walk on the weighted
bipartite graph with biadjacency $\bB$; by Perron--Frobenius, the
multiplicity of $\sigma = 1$ equals the number of connected components,
so $\sigma_2 < 1$ iff the graph is connected, which is also exactly the
condition $V_1 \cap V_2 = \operatorname{span}\{\mathbf{1}\}$
\citep{abowd1999,gaure2013}. \qed

\subsection{Proof of Theorem~\ref{thm:multiG}}

\emph{(i)}  Within one factor the indicator columns have disjoint
supports, so $\mathbf{A}_g'\Omega^\ast\mathbf{A}_g = \bD_g$ (diagonal);
for $g \ne h$, $\mathbf{A}_g'\Omega^\ast\mathbf{A}_h = \bB_{gh}$ by
definition.  Theorem~\ref{thm:rate}(i) identifies the local rate with
$\rho$ of the block Gauss--Seidel operator of $\mathcal{F}$, which is a
function of these blocks only.

\emph{(ii)}  \emph{Eigenspace of $1$.}  $\mathbf{M}v = v \iff
-\mathcal{L}'v = (\mathcal{D} + \mathcal{L})v \iff \mathcal{F}v =
\bzero$, so the fixed points of the sweep are exactly the null vectors
of $\mathcal{F}$; for connected designs these are the $G - 1$
between-factor shift directions, and in general the flats of the
identification quotient.  \emph{Semisimplicity.}  Suppose the eigenvalue
$1$ had a nontrivial Jordan block: then there is $v$ with
$(\mathbf{M} - \bI)v = f$ for some $f \in
\operatorname{null}(\mathcal{F}) \setminus \{\bzero\}$.  Since
$\mathbf{M} - \bI = -(\mathcal{D} + \mathcal{L})^{-1}\mathcal{F}$, this
reads $\mathcal{F}v = -(\mathcal{D} + \mathcal{L})f$.  Multiply on the
left by $f'$: the left side is $f'\mathcal{F}v = (\mathcal{F}f)'v = 0$.
For the right side, $\mathcal{F}f = \bzero$ gives
$(\mathcal{D} + \mathcal{L})f = -\mathcal{L}'f$, hence
$f'\mathcal{D}f + 2f'\mathcal{L}f = f'\mathcal{F}f = 0$, i.e.\
$f'\mathcal{L}f = -\tfrac12 f'\mathcal{D}f$, and therefore
$f'(\mathcal{D} + \mathcal{L})f = \tfrac12 f'\mathcal{D}f > 0$ because
$\mathcal{D}$ is positive definite (every level carries positive
weight) and $f \ne \bzero$.  The two sides disagree---contradiction; so
the eigenvalue $1$ is semisimple and $\R^{\sum J_g} =
\operatorname{null}(\mathcal{F}) \oplus W$ with $W$ an
$\mathbf{M}$-invariant complement on which
$\rho(\mathbf{M}|_W) = \rho_{\mathrm{GS}} < 1$ by
Theorem~\ref{thm:rate}(i) (this is the semiconvergence of Gauss--Seidel
for consistent singular symmetric positive semidefinite systems; cf.\
\citealp{keller1965,berman1994}).  \emph{Differences.}  Decomposing an
arbitrary start $e_0 = f_0 + w_0$, the iterates are $e_t = f_0 +
\mathbf{M}^t w_0$, so $e_{t+1} - e_t = (\mathbf{M} -
\bI)\mathbf{M}^t w_0$, whose norm ratios converge to
$\rho(\mathbf{M}|_W)$ for generic $w_0$: the fixed-point component drops
out automatically, which is what makes plain power iteration on
differences a valid estimator of $\rho_{\mathrm{GS}}$.
\emph{Identification.}  $\rho_{\mathrm{GS}} < 1$ on the quotient iff
$\mathcal{F}$ is positive definite there (Gauss--Seidel on a symmetric
system converges iff the matrix is positive definite on the relevant
space; \citealp[Sec.~11.2]{golub2013}), which is the definition of the
multi-way design being identified up to the flats.  \emph{Cost.}  One
application of $\mathbf{M}$ is one block forward-substitution: for each
$g$ in order, a diagonal solve after accumulating
$\sum_{h \ne g} \bB_{gh} v_h$, i.e., $O\big(\sum_{g<h}
\operatorname{nnz}(\bB_{gh})\big)$ operations.

\emph{(iii)}  Work in the $\Omega^\ast$ inner product on $\R^N$ and let
$\varepsilon_t \in V_1 + \dots + V_G$ be the error of the fitted
fixed-effect vector after sweep $t$.  Updating block $g$ to its exact
blockwise maximizer replaces $\varepsilon$ by $(\bI -
\bP_g)\varepsilon$, where $\bP_g$ is the $\Omega^\ast$-orthogonal
projection onto $V_g$: the block update minimizes the quadratic (IRLS)
energy over corrections in $V_g$, and the minimizer of
$\|\varepsilon + \delta\|_{\Omega^\ast}$ over $\delta \in V_g$ is
$\delta = -\bP_g\varepsilon$.  One sweep therefore applies
$\mathbf{E} = (\bI - \bP_G)\cdots(\bI - \bP_1)$---the multiplicative
subspace-correction (equivalently, cyclic alternating projections onto
the complements $M_g = V_g^{\perp}$)---and the level-space operator
$\mathbf{M}$ of (ii) is the matrix of
$\mathbf{E}|_{V_1 + \dots + V_G}$ in the (redundant) level
parametrization, so the two have the same nonzero spectrum and
$\rho_{\mathrm{GS}} = \rho(\mathbf{E}|_{(\cap_g M_g)^{\perp} \cap (V_1 +
\dots + V_G)}) \le \|\mathbf{E} - \bP_{\cap_g M_g}\|_{\Omega^\ast}$.
The product bound of \citet{smith1977} (see also
\citealp{deutschhundal1997,badea2012}) gives
$\|\mathbf{E} - \bP_{\cap M_g}\|_{\Omega^\ast} \le \big(1 -
\prod_{g=1}^{G-1}(1 - \tilde c_g^{\,2})\big)^{1/2}$ with $\tilde c_g =
c\big(M_g,\, M_{g+1} \cap \dots \cap M_G\big)$; by the complement
duality of the Friedrichs angle, $c(A, B) = c(A^{\perp}, B^{\perp})$
\citep[Ch.~9]{deutsch2001}, and $M_{g+1} \cap \dots \cap M_G =
(V_{g+1} + \dots + V_G)^{\perp}$, so $\tilde c_g = c_g$ as defined in
the statement.  Each $c_g$ is the largest cosine between $V_g \ominus
(\text{intersection})$ and the pooled span of the later factors, i.e.,
the largest non-unit generalized singular value of the level-space
pencil formed from $\bD_g$, the stacked cross-tabulations
$(\bB_{g,g+1}, \dots, \bB_{g,G})$, and the pooled Gram matrix of the
later factors---computable by a small dense or sparse eigenproblem.
Positivity of all angles ($c_g < 1$) is equivalent to $V_g \cap
(V_{g+1} + \dots + V_G) = \text{flats}$ for each $g$, which is
identification.

\emph{(iv)}  For $G = 2$ the sweep operator eigenvalues are those of
$\bD_2^{-1}\bB_{21}\bD_1^{-1}\bB_{12}$ (Appendix~\ref{app:global},
proof of Theorem~\ref{thm:rate}(ii)), whose non-unit spectrum is
$\{\sigma_i^2\}$ of the normalized biadjacency kernel, giving
$\rho_{\mathrm{GS}} = \sigma_2^2 = c_F^2$; the bound in (iii) has a
single factor, $(1 - (1 - c_F^2))^{1/2} = c_F$. \qed

% =========================================================================
\section{Proofs for Section~\ref{sec:asymptotics} of this supplement}\label{app:classical}
% =========================================================================

\subsection{Proof of Theorem~\ref{thm:classical}}

Under Assumption~\ref{ass:classical}, the normalizing regularity and
divergence conditions of \citet[Sec.~3]{fahrmeir1985} hold
for the canonical-link exponential family with bounded covariates:
their normed information $\bF_N^{-1/2}$ exists for large $N$, and
their Theorems~1--2 give existence, consistency, and
$\bF_N^{1/2}(\hat\btheta_N - \btheta_0) \convd \mathcal{N}(\bzero,
\phi \bI)$; Assumption~\ref{ass:classical}(ii) converts this to the
stated $\sqrt N$ form.  Efficiency is the standard exponential-family
Cram\'er--Rao attainment.  Under correct mean but misspecified
distribution, the estimator is the PML estimator of
\citet{gourieroux1984} for the linear exponential family, whose
Theorem~3 gives the sandwich limit; consistency of the cell-accumulated
sandwich estimator follows from the law of large numbers applied to
$\hat u_i^2$-weighted design moments, using the uniform consistency of
$\hat\mu_i$ on bounded designs. \qed

\subsection{Proof of Theorem~\ref{thm:diverging}}\label{app:diverging}

Normalize $\phi = 1$.  Write $\bF = \bF_N(\btheta_0)$ (positive
definite for large $N$ by Assumption~\ref{ass:diverging}(ii), which
presupposes invertibility), $\bU = \bU(\btheta_0)$, and change variables
$\bs = \bF^{1/2}(\btheta - \btheta_0)$, $g(\bs) = \ell(\btheta_0 +
\bF^{-1/2}\bs)$.  Then $g$ is concave, $\nabla g(\bzero) =
\bF^{-1/2}\bU$, and $\nabla^2 g(\bs) = -\bF^{-1/2}\bH(\btheta_{\bs})
\bF^{-1/2}$.

\emph{Step 1: score bound.}  $\E\|\nabla g(\bzero)\|^2 =
\tr(\bF^{-1/2}\,\Var(\bU)\,\bF^{-1/2}) = q_N$ by the information
equality, so $\|\nabla g(\bzero)\| = O_p(\sqrt{q_N})$.

\emph{Step 2: uniform curvature control.}  Fix $r \ge 1$ and let
$\|\bs\| \le r\sqrt{q_N}$.  By Assumption~\ref{ass:diverging}(i,iv),
$\omega_i = b''(\eta_{0i}) \in [\underline\omega, \bar\omega]$ with
$0 < \underline\omega \le \bar\omega < \infty$, and the balanced-leverage
condition gives $\|\bF^{-1/2}\ba_i\|^2 = h_i/\omega_i \le C q_N /
(N\underline\omega)$.  Hence
\[
\max_i |\eta_i(\btheta_{\bs}) - \eta_{0i}|
 = \max_i |\ba_i'\bF^{-1/2}\bs|
 \le r\sqrt{q_N} \cdot \sqrt{C q_N / (N \underline\omega)}
 = C_r\, q_N/\sqrt{N} =: \varepsilon_N \to 0
\]
using $q_N^2/N \to 0$.  With $L$ the Lipschitz constant of $b''$ on the
relevant compact set,
$\big| \sum_i \{b''(\eta_i(\btheta_\bs)) - b''(\eta_{0i})\}
(\ba_i'\bu)^2 \big| \le (L\varepsilon_N/\underline\omega)\, \bu'\bF\bu$
for all $\bu$, so uniformly on the ball,
\begin{equation}\label{eq:hessian-sandwich}
-(1 + \delta_N)\,\bI \preceq \nabla^2 g(\bs) \preceq
-(1 - \delta_N)\,\bI ,
\qquad \delta_N := L\varepsilon_N/\underline\omega \to 0 .
\end{equation}

\emph{Step 3: existence and rate.}  On $\|\bs\| = r\sqrt{q_N}$, a
second-order Taylor expansion with \eqref{eq:hessian-sandwich} gives
$g(\bs) - g(\bzero) \le \|\nabla g(\bzero)\| r\sqrt{q_N} - \tfrac12(1 -
\delta_N) r^2 q_N$.  On the event $\{\|\nabla g(\bzero)\| \le K
\sqrt{q_N}\}$, whose probability exceeds $1 - K^{-2} \E\|\nabla
g(0)\|^2/q_N = 1 - K^{-2}$, the right side is negative for $r > 4K$;
concavity of $g$ then confines the (unique, by strict concavity)
maximizer $\hat\bs$ to the open ball, proving existence and
$\|\hat\bs\| = \|\bF^{1/2}(\hat\btheta - \btheta_0)\| =
O_p(\sqrt{q_N})$.

\emph{Step 4: asymptotic linearity.}  By the integral form of the mean
value theorem, $\bzero = \nabla g(\hat\bs) = \nabla g(\bzero) -
\tilde{\mathbf{J}} \hat\bs$ with $\tilde{\mathbf{J}} = -\int_0^1
\nabla^2 g(t \hat\bs)\dd t$; by \eqref{eq:hessian-sandwich},
$\|\tilde{\mathbf{J}} - \bI\| \le \delta_N$, so
\[
\|\hat\bs - \nabla g(\bzero)\|
 \le \frac{\delta_N}{1 - \delta_N} \|\nabla g(\bzero)\|
 = O_p\!\left( \frac{q_N}{\sqrt N} \sqrt{q_N} \right)
 = O_p\!\left( q_N^{3/2}/\sqrt{N} \right) = o_p(1)
\]
under $q_N^3/N \to 0$.

\emph{Step 5: contrasts.}  For $\bC_N$ with $\bm{\Sigma}_{C,N} =
\bC_N'\bF^{-1}\bC_N \to \bm{\Sigma}_C \succ \bzero$ and any fixed
$\bv \in \R^d$, write $Z_N = \bv'\bm{\Sigma}_{C,N}^{-1/2}\bC_N'
(\hat\btheta - \btheta_0) = \bv'\bm{\Sigma}_{C,N}^{-1/2}\bC_N'
\bF^{-1/2}\hat\bs$.  The remainder from Step 4 contributes
$\le \|\bm{\Sigma}_{C,N}^{-1/2}\bC_N'\bF^{-1/2}\|\, o_p(1) = o_p(1)$
(the matrix has orthonormal-bounded rows by construction).  The main
term is $\sum_i c_{N,i} (y_i - b'(\eta_{0i}))$ with $c_{N,i} =
\bv'\bm{\Sigma}_{C,N}^{-1/2} \bC_N'\bF^{-1}\ba_i$: independent,
mean-zero, total variance $\bv'\bv(1 + o(1))$, and
$\max_i |c_{N,i}| \le \|\bv\|\, \|\bm{\Sigma}_{C,N}^{-1/2}\|\,
\|\bC_N'\bF^{-1/2}\| \max_i \|\bF^{-1/2}\ba_i\| = O(\sqrt{q_N/N}) \to
0$; since the residuals have uniformly bounded $(2+\epsilon)$
moments on compact $\eta$-sets, the Lindeberg condition holds and
$Z_N \convd \mathcal{N}(0, \|\bv\|^2)$; Cram\'er--Wold completes the
claim.  Consistency of the plug-in variance follows from
\eqref{eq:hessian-sandwich} evaluated at $\hat\bs$:
$\|\bF^{-1/2}\bF(\hat\btheta)\bF^{-1/2} - \bI\| \le \delta_N'
\convp 0$, which implies relative consistency of
$\bC_N'\bF(\hat\btheta)^{-1}\bC_N$. \qed

\subsection{Proof of Lemma~\ref{lem:conditional}}

For fixed $\bgamma$ the inner maximization separates over cells;
$\partial \ell/\partial \alpha_j = D_j - e^{\alpha_j} E_j(\bgamma)$ with
$E_j(\bgamma) = \sum_i e^{\bgamma'\bx_{ij} + o_{ij}}$, giving
$e^{\hat\alpha_j} = D_j / E_j$ for $D_j > 0$ (for $D_j = 0$ the supremum
is approached as $\alpha_j \to -\infty$ and the cell contributes a
$\bgamma$-free constant).  Substituting,
\[
\ell_p(\bgamma) = \sum_{j : D_j > 0} \Big\{ \sum_i y_{ij}
(\bgamma'\bx_{ij} + o_{ij}) - D_j \log E_j(\bgamma) \Big\}
 + \sum_j D_j (\log D_j - 1) .
\]
The bracket equals $\sum_i y_{ij}\log \pi_{ij}(\bgamma)$ plus terms free
of $\bgamma$, which is the multinomial log likelihood given $D_j$ up to
the combinatorial constant; and the multinomial law given the total is
the standard Poisson-splitting identity.  Concavity: $\bgamma \mapsto
\log E_j(\bgamma)$ is log-sum-exp of affine maps, hence convex.
Differentiating the bracket gives the stated score and, again, negative
Hessian $\sum_j D_j \Var_{\pi_j}(\bx_j)$; the Schur form follows by
evaluating \eqref{eq:s-moments} at $\hat\balpha(\bgamma)$, where $\mu_{ij}
= D_j \pi_{ij}$, so $M_j^{(0)} = D_j$, $\mathbf{m}_j = D_j
\E_{\pi_j}[\bx_j]$, and $\bQ - \sum_j \mathbf{m}_j\mathbf{m}_j'/M_j^{(0)}
= \sum_j D_j \{\E_{\pi_j}[\bx\bx'] - \E_{\pi_j}[\bx]\E_{\pi_j}[\bx]'\}$.
\qed

\subsection{Proof of Theorem~\ref{thm:proportional}}

Throughout, expectations condition on the (fixed) covariates, offsets,
and cell effects; cells are independent.  Let $\ell_j(\bgamma)$ denote
cell $j$'s conditional (multinomial) log likelihood and $L_J(\bgamma) =
J^{-1}\sum_j \{\ell_j(\bgamma) - \ell_j(\bgamma_0)\}$, so that
$\hat\bgamma_J$ maximizes the concave function $L_J$ by
Lemma~\ref{lem:conditional}.

\emph{(i) Consistency.}  Each summand satisfies, on the compact set
$\Gamma$, $|\ell_j(\bgamma) - \ell_j(\bgamma_0)| \le 2 D_j
\sup_{i,\bgamma\in\Gamma}|(\bgamma-\bgamma_0)'\bx_{ij}| \le C_\Gamma
D_j$, and by Assumption~\ref{ass:proportional}(ii) the cell means are
bounded, so $D_j$ has uniformly bounded moments of all orders; hence
$\Var\{L_J(\bgamma)\} \le C/J \to 0$ and $L_J(\bgamma) - \E L_J(\bgamma)
\convp 0$ pointwise.  The mean is a per-cell Kullback--Leibler
divergence: $\E\{\ell_j(\bgamma) - \ell_j(\bgamma_0)\} = -\E\,
\mathrm{KL}\big(\pi_j(\bgamma_0) \,\|\, \pi_j(\bgamma)\big) \le 0$, and
a second-order expansion at $\bgamma_0$ with
Assumption~\ref{ass:proportional}(iii) gives an $\epsilon_0 > 0$ and,
for large $J$, $\E L_J(\bgamma) \le -\tfrac{\lambda}{4}\|\bgamma -
\bgamma_0\|^2$ whenever $\|\bgamma - \bgamma_0\| \le \epsilon_0$;
concavity of $\E L_J$ (with $\E L_J(\bgamma_0) = 0$) extends the strict
separation to all of $\Gamma$: for $\|\bgamma - \bgamma_0\| \ge
\epsilon$ with $\epsilon \le \epsilon_0$, evaluating along the segment
from $\bgamma_0$ yields $\E L_J(\bgamma) \le
-\tfrac{\lambda\epsilon^2}{4}$.  By the convexity lemma
\citep{andersen1982,pollard1991,niemiro1992}, pointwise convergence of
the concave functions $L_J - \E L_J$ to $0$ is automatically uniform on
compacts, and the argmax of $L_J$ converges in probability to
$\bgamma_0$.

\emph{(ii) Normality.}  The profile score at the truth is
$\bU_J(\bgamma_0) = \sum_j \bs_j$ with $\bs_j = \sum_i (y_{ij} - D_j
\pi_{ij}(\bgamma_0))\bx_{ij}$.  Conditionally on $D_j$ the multinomial
model is correctly specified, so $\E[\bs_j \mid D_j] = \bzero$, hence
$\E \bs_j = \bzero$ \emph{exactly, for every $J$ and every configuration
of cell effects}---this is the absence of incidental-parameter bias.
The summands are independent with $\Var(\bs_j) = \E[D_j
\Var_{\pi_j}(\bx_j)]$ (conditional information equality), so
$\Var(J^{-1/2}\bU_J) = \bar\bF_J(\bgamma_0)$, and $\|\bs_j\| \le
2 D_j \max_i \|\bx_{ij}\|$ has bounded third moments, so Lyapunov's CLT
gives $J^{-1/2}\bU_J \convd \mathcal{N}(\bzero, \bar\bF)$ by
Assumption~\ref{ass:proportional}(iii).  For the Hessian, $J^{-1}\bH_p
(\bgamma) = J^{-1}\sum_j D_j \Var_{\pi_j(\bgamma)}(\bx_j)$ is an average
of independent terms, bounded by $C D_j$ and Lipschitz in $\bgamma$ with
constant $\le C' D_j$ on $\Gamma$; a standard bracketing/finite-cover
argument gives $\sup_{\|\bgamma - \bgamma_0\| \le \epsilon_J}
\|J^{-1}\bH_p(\bgamma) - \bar\bF_J(\bgamma_0)\| \convp 0$ for any
$\epsilon_J \to 0$.  A Taylor expansion of the score around $\bgamma_0$,
solved at $\hat\bgamma_J$, then yields $\sqrt{J}(\hat\bgamma_J -
\bgamma_0) = \bar\bF_J^{-1} J^{-1/2}\bU_J + o_p(1)$ and the stated
limit.

\emph{(iii) Variance estimation.}  Consistency of $J^{-1}
\bH_p(\hat\bgamma_J)$ for $\lim \bar\bF_J$ is the uniform convergence
just shown plus (i).  For the clustered estimator, $J^{-1}\sum_j
\bs_j(\hat\bgamma)\bs_j(\hat\bgamma)'$: at $\bgamma_0$ the LLN with
bounded fourth moments gives convergence to $\lim J^{-1}\sum_j
\Var(\bs_j)$; the difference between evaluation at $\hat\bgamma$ and at
$\bgamma_0$ is bounded by a Lipschitz-in-$\bgamma$ envelope with
integrable constant times $\|\hat\bgamma - \bgamma_0\| \convp 0$.  If
the within-cell joint distribution is misspecified but the conditional
mean $\E[y_{ij}\mid D_j] = D_j\pi_{ij}(\bgamma_0)$ is correct, then
$\E\bs_j = \bzero$ still holds, the same argument applies with
$\Var(\bs_j)$ no longer equal to the conditional information, and the
sandwich $\bH_p^{-1}(\sum_j \bs_j\bs_j')\bH_p^{-1}$ remains consistent;
the $G/(G-1)$ correction is asymptotically negligible.

\emph{(iv) Efficiency.}  In the model with unrestricted
$(\alpha_1, \alpha_2, \dots)$, the nuisance tangent space at
$(\bgamma_0, \balpha_0)$ is spanned, cell by cell, by the
$\alpha_j$-scores $u_j = \sum_i (y_{ij} - \mu_{ij})$.  The efficient
score for $\bgamma$ is the residual of the projection of the
$\bgamma$-score onto this space; by independence across cells the
projection is cellwise, and within cell $j$,
$\E[(\sum_i (y_{ij}-\mu_{ij})\bx_{ij})\, u_j] = \sum_i \mu_{ij}
\bx_{ij}$, $\E[u_j^2] = \sum_i \mu_{ij}$, so the projection residual is
$\sum_i (y_{ij} - \mu_{ij})\bx_{ij} - \frac{\sum_i \mu_{ij}
\bx_{ij}}{\sum_i \mu_{ij}} \sum_i (y_{ij} - \mu_{ij}) = \bs_j$ evaluated
at the truth: exactly the conditional score.  The asymptotic variance in
(ii) equals the inverse of the efficient information, so the bound is
attained \citep[Ch.~25]{vandervaart1998}. \qed

% =========================================================================
\section{Proofs for Section~\ref{sec:extensions} of this supplement}\label{app:extensions}
% =========================================================================

\subsection{Proof of Proposition~\ref{prop:penalized}}

(i) Adding $\tfrac{\lambda_2}{2}\|\bgamma\|^2$ contributes $\lambda_2
\bI_p$ to the $\bgamma$-block of the curvature and $-\lambda_2\bgamma$
to its gradient; the modified system has the same dimensions and cost.
(ii) The proximal Newton step minimizes the penalized quadratic model
$m(\btheta) = -\bU'(\btheta - \btheta_t) + \tfrac12 (\btheta -
\btheta_t)'\bF(\btheta - \btheta_t) + \lambda_1\|\bgamma\|_1$ (plus the
ridge term); one coordinate update needs the corresponding row of $\bF$
and the current residual of the linear system, both dense in
$\R^{p+k}$, so a full cycle costs $O((p+k)^2)$---all references to data
enter through $(\bU, \bF)$, which Theorem~1 of the main paper provides from
one pass.  Since $-\ell$ is convex with locally Lipschitz Hessian and
$\bF \succ \bzero$ near the solution, the inexact proximal Newton method
converges locally quadratically \citep[Thms.~3.4--3.5]{leesun2014}.
(iii) The KKT conditions are $|\,[\bU]_r - [\bF(\btheta -
\btheta_t)]_r\,| \le \lambda_1$ for inactive $\gamma_r$, with equality
sign structure for active ones: functions of $(\bU, \bF)$ only. \qed

\subsection{Proof of Proposition~\ref{prop:firth}}

For canonical links, $\bF(\btheta) = \sum_i b''(\eta_i)\ba_i\ba_i'$
depends on $\btheta$ through the $\eta_i$, and $\partial \bF /
\partial\theta_r = \sum_i b'''(\eta_i) a_{ir}\, \ba_i\ba_i'$.  Firth's
adjustment \citep{firth1993} is $A_r = \tfrac12 \tr(\bF^{-1}
\partial\bF/\partial\theta_r) = \tfrac12 \sum_i b'''(\eta_i) a_{ir}\,
\ba_i'\bF^{-1}\ba_i = \sum_i h_i\, \frac{b'''(\eta_i)}{2 b''(\eta_i)}
a_{ir}$ with $h_i = b''(\eta_i)\ba_i'\bF^{-1}\ba_i$: the adjusted score
is the ordinary score with working residual $y_i - b'(\eta_i) + h_i
b'''(\eta_i)/2b''(\eta_i)$, which the cell accumulation of
Theorem~1 of the main paper ingests unchanged.  For the leverages, write
the partitioned inverse blocks $\bV_{\gamma\gamma}, \bV_{\gamma\delta},
\bV_{\delta\delta}$ of $\bF^{-1}$ and expand $\ba_i =
(\bx_i', \bw_{j(i)}')'$:
$\ba_i'\bF^{-1}\ba_i = \bx_i'\bV_{\gamma\gamma}\bx_i + 2\bx_i'
(\bV_{\gamma\delta}\bw_{j(i)}) + \bw_{j(i)}'\bV_{\delta\delta}
\bw_{j(i)}$.  Precomputing $\bc_j = \bV_{\gamma\delta}\bw_j$ ($O(Jkp)$)
and $d_j = \bw_j'\bV_{\delta\delta}\bw_j$ ($O(Jk^2)$) leaves per-datum
cost $O(p^2 + p)$, giving the stated totals; $b'''/2b''$ equals $1/2$
for Poisson and $(1-2\mu_i)/2$ for the logit by direct computation.
The quasi-scoring iteration and its properties are those of
\citet{kosmidis2009}; finiteness of the penalized maximizer for binomial
responses is Theorem~1 of \citet{heinze2002}. \qed


\bibliographystyle{chicago}
\bibliography{references}

\end{document}
